\documentclass{article}
\usepackage[dblblindworkshop]{neurips_2026} % After being accepted, the authors should add "final":  % \usepackage[dblblindworkshop, final]{neurips_2026}

\usepackage[utf8]{inputenc} % allow utf-8 input
\usepackage[T1]{fontenc}    % use 8-bit T1 fonts
\usepackage{hyperref}       % hyperlinks
\usepackage{url}            % simple URL typesetting
\usepackage{booktabs}       % professional-quality tables
\usepackage{amsfonts}       % blackboard math symbols
\usepackage{nicefrac}       % compact symbols for 1/2, etc.
\usepackage{microtype}      % microtypography
\usepackage{xcolor}         % colors

\usepackage{amsmath}
\usepackage{graphicx}

\title{Depth Compression in Time-Series Foundation Models: Saturation, Recovery, and Measured Speedups in Chronos-2}
% or Compressing Chronos-2 by Depth: Saturation, Latency, and Dataset-Dependent Truncation
\workshoptitle{TS-LIMITS @NeurIPS 2026 -- Generalization for Time Series in Tight Settings: Latency, Inference, Memory, prIvacy and susTainability}

\author{Anonymous Authors}

\begin{document}


\maketitle


\begin{abstract}
Time-series foundation models (TSFMs) support zero-shot forecasting but can be too expensive for latency- and memory- constrained deployments. Most compression methods reduce parameter count or numerical precision. In this work, we ask whether model depth itself can be reduced by skipping later layers. Using Chronos-2, we find that forecasting performance using a direct linear readout saturates before the final representation under a 5\% criterion on all tested datasets, although the saturation depth varies across datasets. Effective rank and CKA reveal clear layerwise structure but do not identify a universal transition associated with saturation. We then attach lightweight linear adapters to intermediate representations and reuse the frozen Chronos-2 quantile head. Dataset-specific adapters often recover much of the native forecasting performance, whereas cross-dataset transfer is less consistent. We show that truncating at layer 3 drops \(71\%\) of the active parameters, reduces peak GPU memory by \(34\%\) (on batches of \(256\)), and provides a \(\times 3.23\) measured speedup, with a MASE increase of \(7\%\) averaged across three datasets. Depth compression can therefore provide substantial efficiency gains, with the optimal truncation depth and associated gains varying by dataset.
% but no saturation- or geometry-based criterion determines a universal truncation depth.
\end{abstract}


\section{Introduction} \label{intro}

Time-series foundation models (TSFMs) provide strong zero-shot forecasts, but their inference cost can prevent their use in latency- and memory-constrained settings \citep{ansari2024chronoslearninglanguagetime, das2024decoderonlyfoundationmodeltimeseries, woo2024unifiedtraininguniversaltime, ansari2025chronos2univariateuniversalforecasting}. Production monitoring and control systems may forecast thousands of series on a fixed schedule, often under limited computational resources. In these settings, latency and memory matter alongside accuracy.

Common compression methods reduce numerical precision \citep{frantar2023gptqaccurateposttrainingquantization} or distill a model into a smaller architecture \citep{hinton2015distillingknowledgeneuralnetwork, sanh2020distilbertdistilledversionbert}. These approaches do not directly address whether later layers can be removed. Studies of language models have found that later transformer layers can contribute little to the final prediction, motivating structured layer dropping \citep{fan2019reducingtransformerdepthdemand}, layer pruning \citep{gromov2025unreasonableineffectivenessdeeperlayers, men2024shortgptlayerslargelanguage}, and early exit \citep{teerapittayanon2017branchynetfastinferenceearly, schwartz2020righttooljobmatching, Elhoushi_2024}. Recent work has also found that intermediate TSFM layers can outperform the final layer as frozen features for time-series classification \citep{feofanov2026mantisv2closingzeroshotgap}. It remains unclear whether similar depth redundancy exists for forecasting and, if so, how reliably a usable truncation depth can be selected.

To test how much forecasting information is available at intermediate depths, we ask the following questions. First, does forecasting performance under a simple linear readout saturate across depth in Chronos-2 \citep{ansari2025chronos2univariateuniversalforecasting}? If so, is this saturation reflected in the representation geometry, independently of the probe? Finally, can intermediate representations recover near-full-model forecasting performance? We first test whether a dataset-specific linear adapter can align intermediate representations with the pretrained forecasting head, then examine whether this alignment transfers across datasets. We additionally compare truncation depths selected from forecasting saturation and representation geometry and quantify the resulting compression-performance trade-offs.

\paragraph{Contribution.}
\begin{enumerate}
    \item \textbf{Layerwise forecasting saturation.}
    We characterize forecasting performance across Chronos-2 depth on datasets from and outside the documented pretraining corpus and analyze how saturation relates to representation geometry.

    \item \textbf{Recovery from intermediate representations.}
    We test whether lightweight linear adapters can approach native forecasting performance from intermediate representations and whether the learned alignment transfers across datasets.

    \item \textbf{Measured depth compression.}
    We construct physically truncated models and directly measure latency and memory, showing that depth compression can provide substantial efficiency gains but that the useful truncation depth remains dataset-dependent.
\end{enumerate}


\section{Experimental Framework}

\subsection{Models and forecasting setup} \label{sec:model-setup}

\paragraph{Model and layer extraction.} We use Chronos-2 \citep{ansari2025chronos2univariateuniversalforecasting}, a ${\sim}120$M-parameter encoder-only time-series foundation model with 12 transformer encoder blocks. All Chronos-2 parameters are frozen. We use a context length \(C=512\), forecast horizon \(H=64\), and patch size \(P=16\), which corresponds to \(K=4\) forecast slots. Each input time series follows the unmodified Chronos-2 processing pipeline up to the representation being tested. We extract representations at 14 points throughout the network: the token embeddings before the first encoder block (\(\mathrm{Emb}\)), the outputs of the 12 encoder blocks (\(L1 - L12\)), and the post-RMSNorm output after the last encoder block and before the Chronos-2 quantile head, denoted \(L12 + \mathrm{RMS}\).

\paragraph{Shared forecast-token probe.} At each representation point, a single linear probe is shared across all \(K\) forecast slots. Each representation dimension is standardized using statistics from the training forecast-slot states at that layer. The probe then maps each slot's 768-dimensional hidden state to \(Q\times P\) outputs, where \(Q\) is the number of quantiles. 
% The resulting \(K\) predicted patches are concatenated along the forecast horizon to produce the full \(Q\times H\) forecast.

\paragraph{Truncated models.}
At each depth \(\ell\), let \(h_\ell\) denote the forecast-slot representation after encoder block \(\ell\). We construct a \textit{truncated model} by stopping the Chronos-2 encoder after block \(\ell\) and attaching a layer-specific linear adapter
\(A_\ell:\mathbb{R}^{768}\rightarrow\mathbb{R}^{768}\), shared across the \(K\) forecast slots.
The adapted representations are passed through the frozen final RMSNorm and pretrained forecasting head,
\[
h_\ell \rightarrow A_\ell(h_\ell)
\rightarrow \mathrm{RMSNorm}_{\mathrm{final}}
\rightarrow \mathrm{head}.
\]
All Chronos-2 parameters remain frozen, and only \(A_\ell\) is trained. When an adapter is trained and evaluated on the same dataset, we refer to it as a \textit{dataset-specific linear adapter}.

\subsection{Data and evaluation} \label{sec:data-eval}

\paragraph{ID/OOD protocol and datasets.}
We use seven datasets, all evaluated at hourly resolution. Four datasets are listed in the documented Chronos-2 pretraining corpus \citep{ansari2025chronos2univariateuniversalforecasting} and are referred to as \textit{PT-ID}: M4 \citep{MAKRIDAKIS202054}, Electricity \citep{godahewa2021monashtimeseriesforecasting}, Uber TLC \citep{fivethirtyeight_uber_tlc}, and Wind Farms \citep{godahewa2021monashtimeseriesforecasting}. The remaining three datasets are not listed in the pretraining corpus and are referred to as \textit{PT-OOD}: BOOM \citep{cohen2025timedifferentobservabilityperspective}, SG Carpark \citep{qiao2026itstimegenerationtime}, and Coastal T-S \citep{qiao2026itstimegenerationtime}. Dataset descriptions are provided in Appendix~\ref{apndx:datasets}, with split construction, window selection, and filtering details in Appendix~\ref{apndx:data-splits}.

\paragraph{Forecasting metrics.}
For the saturation analysis, we train and evaluate the probes using the original Chronos-2 quantile loss \citep{ansari2025chronos2univariateuniversalforecasting}. To compare truncated models with native Chronos-2, we use mean absolute scaled error (MASE) \citep{HYNDMAN2006679}, normalized by the seasonal-naive error over the \(C=512\) context window with seasonal period \(m=24\). Probe results are averaged over three probe initializations. Each adapter is fit once. We estimate 95\% confidence intervals using 5,000 bootstrap resamples. Bootstrap details are provided in Appendix~\ref{app:implementation}.


%%%%%%%%%%%%%%%%%% RESULTS %%%%%%%%%%%%%%%%%
\section{Results}

\subsection{Does linearly decodable forecasting performance saturate across depth?}
\label{sec:tunnel}

\paragraph{Saturation region.}
We define the entrance to the \textit{saturation region} as the first layer whose validation loss is within \(\varepsilon\%\) of the final-layer loss. We use \(\varepsilon=5\) in the main analysis and report additional results for \(\varepsilon=2\) in Appendix~\ref{apndx:pt-id-results}. This follows the ``tunnel entrance'' criterion from the representation-learning literature \citep{masarczyk2023tunneleffectbuildingdata}.

The top panels of Figure~\ref{fig:id_tunnel} show the layerwise test quantile loss for two representative PT-ID datasets and one PT-OOD dataset. For Electricity, test loss seems to change little after \(L3\), but the validation curve continues to improve and does not enter the \(5\%\) saturation region until \(L10\). M4 saturates much earlier at \(L3\), after which test performance remains nearly unchanged. SG Carpark also shows a sharp early reduction followed by a visually flatter profile, but its intermediate-layer losses remain above the final representation, and its validation-based saturation region begins only at \(L11\). Results for all PT-ID and PT-OOD datasets under both \(\varepsilon=5\%\) and \(\varepsilon=2\%\) criteria are provided in Appendix~\ref{apndx:pt-id-results}.

Across all tested datasets, linearly decodable forecasting performance saturates before the final representation, both within and outside the documented Chronos-2 pretraining corpus, and the saturation depth is dataset-dependent. We next ask whether this behavior is reflected in the representation geometry, independently of the probe.

\begin{figure}[t]
    \centering
    \includegraphics[width=0.95\linewidth]{figures/main_id_testloss_erank_2x2.png}
    \caption{
    Layerwise forecasting performance and representation geometry for Electricity and M4 (PT-ID) and SG Carpark (PT-OOD). \textbf{Top:} mean test quantile loss across three probe initializations using a single quantile (\(Q=1\)). Shaded bands show 95\% bootstrap confidence intervals. The saturation point (green) is selected independently from the validation curve using the 5\% criterion. \textbf{Bottom:} effective rank of the forecast-slot representations computed on the training set. Stars mark the maximum effective rank. Results for the other PT-ID and PT-OOD datasets are in Appendix \ref{apndx:pt-id-results}.
    }
    \label{fig:id_tunnel}
\end{figure}

\subsection{Is forecasting saturation reflected in representation geometry?}
\label{sec:geometry}

\paragraph{Effective rank.} To characterize the dimensional structure of the representations, we use \textit{effective rank} \citep{7098875}. See Appendix \ref{apndx:effective-rank} for details.

The bottom panels of Figure~\ref{fig:id_tunnel} show the effective-rank profiles of the forecast-slot representations across depth. Effective rank increases from the embedding representation, peaks at \(L6\) for all three datasets, and then declines toward the final layers. However, this peak does not consistently align with the validation-selected saturation point. For Electricity and SG Carpark, the validation losses at \(L6\) are still \(12.5\%\) and \(25.7\%\) above the final-layer loss, respectively. For M4, effective rank peaks inside the saturation region. Similar trends appear for the remaining PT-ID and PT-OOD datasets in Appendix~\ref{apndx:geometry}. Additional linear centered kernel alignment (CKA) results in Appendix~\ref{apndx:cka} also show clear layerwise structure, but its relation to forecasting saturation varies across datasets.

Representation analysis reveals substantial layerwise structure, although its location varies across datasets and does not consistently align with the forecasting saturation point. We next explore whether intermediate representations can support compressed forecasting with performance comparable to the full model.

\subsection{Can intermediate representations recover full-model performance?}
\label{sec:recover-chronos-performance}

\subsubsection{Can a dataset-specific adapter recover native performance?}
\label{sec:dataset-specific-adapter}

\paragraph{Comparison conditions.}
At each intermediate layer, we compare the dataset-specific adapter introduced in Section~\ref{sec:model-setup} against three reference conditions. \textit{Native Chronos-2} denotes the full frozen model. The \textit{frozen native head} passes the layer representation directly through the frozen final RMSNorm and pretrained forecasting head. The \textit{fresh linear probe} is the same layer-specific shared forecast-slot probe used for the saturation analysis in Section~\ref{sec:tunnel}.

In Figure~\ref{fig:dataset-specific-adapter}, the frozen native head shows a substantial gap relative to native Chronos-2 across all three datasets. The fresh linear probe performs considerably better, but remains farther from native Chronos-2 at earlier layers. In contrast, the dataset-specific adapter approaches native Chronos-2 performance much earlier. Results for the remaining datasets show a similar pattern and are provided in Appendix~\ref{apndx:adapter-recovery}. We next ask whether this learned alignment transfers across datasets.

\begin{figure*}[t]
    \centering
    \includegraphics[width=0.95\linewidth]{figures/main_nha_mase__mase.png}
    \caption{
    Dataset-specific recovery of native Chronos-2 forecasting performance for Electricity and M4 (PT-ID) and SG Carpark (PT-OOD). We compare the dataset-specific linear adapter (\textit{Linear adapter + frozen native head}) with three reference conditions: native Chronos-2, the frozen native head applied directly to each intermediate representation after the frozen final RMSNorm, and the fresh linear probe used in the saturation analysis in Section~\ref{sec:tunnel}. Shaded bands show 95\% bootstrap confidence intervals. \(L12 + \mathrm{RMS}\) is omitted.
    }
    \label{fig:dataset-specific-adapter}
\end{figure*}

\subsubsection{Does the recovered alignment transfer across datasets?}
\label{sec:transfer-adapter}

We freeze linear adapters trained on Electricity or M4 and evaluate them without refitting on one PT-ID and two PT-OOD datasets. We compare each transferred adapter with native Chronos-2 and a target-specific adapter.

Figure~\ref{fig:adapter-transfer} shows that transfer is strongly dataset-dependent. Between Electricity and M4, transferred adapters follow the target-specific trend but remain worse across most truncation points. On SG Carpark, both source adapters transfer more effectively and approach the target-specific adapter in later layers. For Coastal T-S, estimates are less precise. Thus, the recovered alignment can transfer, but not consistently across datasets. Full transfer results are provided in Appendix~\ref{apndx:adapter-transfer}.

\begin{figure*}[t]
    \centering
    \includegraphics[width=0.95\linewidth]{figures/main_nha_transfer__mase.png}
    \caption{
    Transfer of adapters trained on Electricity (top) and M4 (bottom) to one PT-ID and two PT-OOD datasets. Transferred adapters are evaluated without refitting and compared with native Chronos-2 and a target-specific adapter trained directly on each target dataset. Shaded bands show 95\% bootstrap confidence intervals. \(L12 + \mathrm{RMS}\) is omitted.
    }
    \label{fig:adapter-transfer}
\end{figure*}


\subsection{Compression-performance trade-offs: How should the truncation depth be selected?}
\label{sec:model-setup-cost}

\paragraph{Depth selection and compression metrics.} We consider two truncation depths: the entrance to the saturation region selected from validation loss and the effective-rank peak computed from the training representations. For each truncated model, we report the fraction of active parameters, measured inference speedup, and relative change in test MASE, all with respect to native Chronos-2. Speedup is the median prediction latency of native Chronos-2 divided by that of the truncated model at batch size 256. Full latency, throughput, and memory measurements are provided in Appendix~\ref{app:latency-methodology}.

Table~\ref{tab:compression-selection} shows that the compression-performance trade-off achieved with dataset-specific adapters varies substantially across datasets. For Uber TLC, Wind Farms, and BOOM, the saturation criterion selects \(L3\), which retains \(29.4\%\) of the active parameters and provides a \(\times 3.23\) speedup. The corresponding MASE increases are \(8.6\%\), \(1.5\%\), and \(10.7\%\), respectively. In contrast, Electricity and SG Carpark saturate much later, at \(L10\) and \(L11\), and therefore provide smaller speedups. M4 saturates at \(L3\), but truncation at this depth increases MASE by \(55.3\%\). Moving to the effective-rank peak at \(L6\) reduces this gap to \(33.8\%\), while still providing a \(\times 1.85\) speedup.

The two depth selection criteria can favor different balances between forecasting performance and compression depending on the dataset. Neither criterion defines a universal truncation point, and the preferred depth depends on the desired balance between forecasting performance and inference cost.

\begin{table*}[t]
    \centering
    \scriptsize
    \setlength{\tabcolsep}{2.5pt}
    \renewcommand{\arraystretch}{0.95}

    \caption{
    Compression-performance trade-off under saturation- and effective-rank-based truncation.
    Speedup is measured at batch size 256 relative to native Chronos-2.
    \(\Delta\)MASE is the relative test MASE increase over native Chronos-2 with 95\% paired bootstrap confidence intervals.
    }
    \label{tab:compression-selection}

    \begin{tabular}{lcccccccc}
    \toprule
    & \multicolumn{4}{c}{Saturation-selected}
    & \multicolumn{4}{c}{Effective-rank-selected} \\
    \cmidrule(lr){2-5}
    \cmidrule(lr){6-9}
    Dataset
    & Depth & Act. params & Speedup & \(\Delta\)MASE [95\% CI]
    & Depth & Act. params & Speedup & \(\Delta\)MASE [95\% CI] \\
    \midrule

    Electricity
    & \(L10\) & \(84.7\%\) & \(\times 1.18\) & \(+14.6\%\,[11.6,17.8]\)
    & \(L6\) & \(53.1\%\) & \(\times 1.85\) & \(+13.7\%\,[10.5,17.5]\) \\

    Uber TLC
    & \(L3\) & \(29.4\%\) & \(\times 3.23\) & \(+8.6\%\,[7.2,10.1]\)
    & \(L6\) & \(53.1\%\) & \(\times 1.85\) & \(+3.9\%\,[3.0,4.7]\) \\

    M4
    & \(L3\) & \(29.4\%\) & \(\times 3.23\) & \(+55.3\%\,[45.0,67.0]\)
    & \(L6\) & \(53.1\%\) & \(\times 1.85\) & \(+33.8\%\,[26.2,42.6]\) \\

    Wind Farms
    & \(L3\) & \(29.4\%\) & \(\times 3.23\) & \(+1.5\%\,[0.1,12.7]\)
    & \(L8\) & \(68.9\%\) & \(\times 1.44\) & \(+3.3\%\,[1.1,20.9]\) \\

    \midrule

    SG Carpark
    & \(L11\) & \(92.6\%\) & \(\times 1.08\) & \(+6.5\%\,[3.4,10.2]\)
    & \(L6\) & \(53.1\%\) & \(\times 1.85\) & \(+12.6\%\,[8.9,16.8]\) \\

    Coastal T-S
    & \(L1\) & \(13.6\%\) & \(\times 6.42\) & \(+1.8\%\,[{-8.5},13.5]\)
    & \(L8\) & \(68.9\%\) & \(\times 1.44\) & \(0.0\%\,[{-8.3},9.5]\) \\

    BOOM
    & \(L3\) & \(29.4\%\) & \(\times 3.23\) & \(+10.7\%\,[6.0,19.6]\)
    & \(L8\) & \(68.9\%\) & \(\times 1.44\) & \(+15.0\%\,[8.4,26.9]\) \\

    \bottomrule
\end{tabular}
\end{table*}

\section{Conclusion and Future Work}

Forecasting performance under a linear readout saturates before the final representation under our default \(\varepsilon=5\%\) criterion on all tested datasets. Representation geometry shows a related but distinct layerwise structure, and intermediate representations can often recover much of the native forecasting performance through lightweight dataset-specific adapters. Chronos-2 can often be truncated while reusing its quantile head: truncation at \(L3\) retains \(29.4\%\) of the active parameters and provides a \(\times 3.23\) measured speedup, with an average MASE increase of \(7\%\) across three datasets. However, no saturation- or geometry-based criterion determines a universal truncation depth.

Future work should test whether these findings generalize across frequencies, horizons, and TSFM families \citep{ansari2024chronoslearninglanguagetime, das2024decoderonlyfoundationmodeltimeseries, woo2024unifiedtraininguniversaltime}. It could also combine depth truncation with quantization or width pruning and consider representation measures beyond effective rank and CKA.


\newpage
%%%%%%%%%%%%%%%% REFERENCES %%%%%%%%%%%%%%
\bibliographystyle{plainnat}
\bibliography{references}

%%%%%%%%%%%%%%%%%%%%%%%%%%%%%%%%%%%%%%%%%%%%%%%%%%%%%%%%%%%%
\newpage
\appendix

\section{Additional Results}

% \subsection{Content-Pooled Representation} \label{apndx:content-pooled-probe}

% % {\color{red} [TODO, Zibo]} setup + saturation with ridge R2 and 5-quantile-bin classification \\

% This appendix asks whether the saturation in
% Section~\ref{sec:tunnel} is specific to the shared forecast-slot
% probe. We repeat the layerwise analysis with linear probes on the
% mean of the content (non-padding) patch states. The probe closest to
% the main text is a quantile probe predicting the full $H$-step
% forecast (21 quantiles); we also report a ridge regression and a
% 5-quantile-bin logistic classifier (bins from the training split),
% which predict the context-normalized future mean, a coarser summary
% of the horizon. Regularization for the quantile and ridge probes is
% selected per layer on an $80/20$ split of the training windows; each
% probe is refit on full train and scored on test.

% This is a robustness check rather than a controlled ablation: the
% pipeline is separate from the main text, with a single seed and no
% held-out validation split. Electricity uses a within-series temporal
% cut; M4's series are too short for that, so M4 falls back to a
% disjoint $70/30$ split over its 414 series, and the two datasets are
% compared in shape, not level. Representations were extracted without
% the embedding and post-RMSNorm taps, so all values are relative to
% $L12$ (counts over $L1$--$L12$). Because the reference differs from
% the main text and there is no validation split, the $5\%$ band is a
% flatness measure, not the tunnel-entrance criterion.

% A flat region is present under this readout as well
% (Figure~\ref{fig:content-flatness}). Under the quantile loss
% (Electricity), 10 of 12 layers lie within $5\%$ of $L12$, the
% exceptions being $L3$ ($5.2\%$ better) and $L9$ ($5.5\%$ worse). The
% coarser targets saturate too: under binned accuracy, 8 of 12 layers
% are in band on each dataset, with excursions mostly upward
% ($L6$--$L8$ beat the final layer on Electricity, peak $L7$ at
% $+6.5\%$; on M4 every layer from $L2$ on matches or beats $L12$, peak
% $L8$ at $+8.0\%$); ridge $R^2$ (Electricity) is the least flat, with
% 5 of 12 in band and $L8$ missing it by $0.1$ percentage points. On
% Electricity all three readouts put their worst layer at $L9$; this
% rules out seed effects but not sampling variation, and it does not
% replicate on M4 ($L9$ sits $6.2\%$ above $L12$), so we do not
% interpret it. All numbers are single-seed test-set point estimates
% without confidence intervals and are read as descriptive; within
% those limits, the mid-network saturation of
% Section~\ref{sec:tunnel} is also seen under a content-pooled readout.

% \begin{figure}[h]
%     \centering
%     \includegraphics[width=0.85\linewidth]{figures/appendix_content_pooling_flatness.png}
%     \caption{
%     Content-mean-pooling performance on Electricity, relative to the
%     last encoder block ($L12$), for the three probes. All curves are
%     oriented so that above $100\%$ is better than $L12$; the quantile
%     loss is mirrored about $100$ for this reason. Shaded: $\pm5\%$
%     band. Ridge $R^2$ at $L1$ and $L2$ falls below the axis range and
%     is annotated. M4 is not shown; its binned-accuracy values are
%     reported in the text.
%     }
%     \label{fig:content-flatness}
% \end{figure}



\subsection{Additional Saturation Results}
\label{apndx:pt-id-results}

Table~\ref{tab:saturation-points} summarizes the validation-selected saturation points across all seven datasets for \(\varepsilon=5\%\) and the stricter \(\varepsilon=2\%\) criterion. Figure~\ref{fig:additional-layerwise-results} shows the corresponding layerwise forecasting performance and effective rank across all datasets.

\begin{table}[t]
    \centering
    \caption{
    Validation-selected saturation points across all datasets under the
    \(\varepsilon=5\%\) criterion used in the main analysis and the stricter
    \(\varepsilon=2\%\) criterion.
    }
    \label{tab:saturation-points}
    \begin{tabular}{lccc}
        \toprule
        Dataset & Pretraining status & \(\varepsilon=5\%\) & \(\varepsilon=2\%\) \\
        \midrule
        Electricity & PT-ID & \(L10\) & \(L12+\mathrm{RMS}\) \\
        M4 & PT-ID & \(L3\) & \(L6\) \\
        Uber TLC & PT-ID & \(L3\) & \(L4\) \\
        Wind Farms & PT-ID & \(L3\) & \(L4\) \\
        \midrule
        SG Carpark & PT-OOD & \(L11\) & \(L12\) \\
        BOOM & PT-OOD & \(L3\) & \(L5\) \\
        Coastal T-S & PT-OOD & \(L1\) & \(L1\) \\
        \bottomrule
    \end{tabular}
\end{table}

\begin{figure*}[t]
    \centering
    \includegraphics[width=\linewidth]{figures/appendix_id_saturation_sensitivity_all7.png}
    \caption{
    Layerwise forecasting performance and representation geometry across all seven datasets.
    Top panels show mean test quantile loss across three probe initializations using a single quantile (\(Q=1\)), with shaded bands indicating 95\% bootstrap confidence intervals.
    Solid and dashed green lines mark the saturation entrances selected independently from the validation curves using \(\varepsilon=5\%\) and \(\varepsilon=2\%\), respectively.
    Bottom panels show effective rank computed from the training forecast-slot representations, with stars marking the maximum effective rank.
    }
    \label{fig:additional-layerwise-results}
\end{figure*}

\subsection{Additional Representation Geometry Results}
\label{apndx:geometry}

The bottom panels of Figure \ref{fig:additional-layerwise-results} show the effective-rank profiles for Uber TLC, Wind Farms, BOOM, and Coastal T-S. Table \ref{tab:geometry-summary} summarizes the relation between the effective-rank peak and forecasting saturation across all seven datasets.


\begin{figure*}[t]
    \centering
    \includegraphics[width=0.85\linewidth]{figures/appendix_id_cka.png}
    \caption{
    Linear CKA between forecast-slot representations across the 14 representation points for the PT-ID and PT-OOD datasets. Higher values indicate greater representational similarity. Dashed lines mark the validation-selected 5\% saturation points.
    }
    \label{fig:additional-cka}
\end{figure*}

\begin{table*}[t]
    \centering
    \caption{
    Summary of effective-rank peaks and their relation to forecasting saturation across all datasets. ``Gap vs. final'' reports the relative validation-loss difference at the effective-rank peak compared with the final \(L12+\mathrm{RMS}\) representation, with positive values indicating worse forecasting performance. The saturation point is selected using the 5\% validation criterion.
    }
    \label{tab:geometry-summary}
    \begin{tabular}{lcccc}
        \toprule
        Dataset & Status & ER peak & Gap vs. final & Saturation point \\
        \midrule
        Electricity & PT-ID  & \(L6\) & \(+12.5\%\) & \(L10\) \\
        Uber TLC    & PT-ID  & \(L6\) & \(-1.1\%\)  & \(L3\) \\
        M4          & PT-ID  & \(L6\) & \(+1.8\%\)  & \(L3\) \\
        Wind Farms  & PT-ID  & \(L8\) & \(-2.2\%\)  & \(L3\) \\
        \midrule
        SG Carpark  & PT-OOD & \(L6\) & \(+25.7\%\) & \(L11\) \\
        Coastal T-S & PT-OOD & \(L8\) & \(-0.6\%\)  & \(L1\) \\
        BOOM        & PT-OOD & \(L8\) & \(+4.9\%\)  & \(L3\) \\
        \bottomrule
    \end{tabular}
\end{table*}

\subsection{Additional Compressed-Model Results}
\label{apndx:compressed-models}

\subsubsection{Dataset-Specific Adapter Recovery}
\label{apndx:adapter-recovery}

Figure \ref{fig:appendix-adapter-recovery} shows the dataset-specific adapter results for all PT-ID and PT-OOD datasets. Across datasets, the linear adapter generally reduces the gap between intermediate representations and the frozen native forecasting head, although the amount of recovery varies with dataset and truncation depth.

\begin{figure*}[t]
    \centering
    \includegraphics[width=\linewidth]{figures/appendix_nha_mase__mase.png}
    \caption{
    Dataset-specific adapter recovery across all PT-ID and PT-OOD datasets.
    }
    \label{fig:appendix-adapter-recovery}
\end{figure*}

\subsubsection{Cross-Dataset Adapter Transfer}
\label{apndx:adapter-transfer}

Figure \ref{fig:appendix-adapter-transfer} shows the complete transfer results for adapters trained on Electricity and M4. Each adapter is frozen and evaluated without refitting on the remaining PT-ID and PT-OOD datasets. The target-specific adapter is shown as a reference for the performance obtained when the adapter is trained directly on the target dataset.

\begin{figure*}[t]
    \centering
    \includegraphics[width=\linewidth]{figures/appendix_nha_transfer__mase.png}
    \caption{
    Cross-dataset transfer of adapters trained on Electricity (top two rows) and M4 (bottom two rows). Each transferred adapter is evaluated without refitting on the remaining PT-ID and PT-OOD datasets. Target-specific adapters are shown as a reference. Dashed lines mark native Chronos-2 performance, and shaded bands show 95\% bootstrap confidence intervals.
    }
    \label{fig:appendix-adapter-transfer}
\end{figure*}

\section{Representation Diagnostics}

\subsection{Effective Rank} \label{apndx:effective-rank}

We use effective rank \citep{7098875} to quantify how broadly representation variance is distributed across feature directions.

First, we center a representation matrix \(X\in\mathbb{R}^{N\times d}\), where \(N\) is the number of sampled forecast-slot states and \(d=768\) is the hidden-state dimension:
\[
X_c = X-\bar{X}.
\]

Then, we compute the singular values \(s_i\) of \(X_c\) and define the normalized variance spectrum as
\[
p_i=\frac{s_i^2}{\sum_j s_j^2}.
\]

The effective rank of the matrix \(X_c\) is defined as
\[
r_{\mathrm{eff}}
=
\exp\left(-\sum_i p_i\log p_i\right).
\]

For intuition, \(r_{\mathrm{eff}}=1\) when all variance lies along one direction. If it is distributed equally across \(k\) directions, then \(r_{\mathrm{eff}}=k\).

For each representation point, the training split contains \(n=1394\) forecasting windows, with \(K=4\) forecast-slot states of dimension \(d=768\) per window. We reshape the representations as
\[
(n,K,d)\rightarrow(nK,d),
\]
which gives \(5576\times768\) available forecast-slot states. We sample \(N=4096\) rows without replacement with seed 0 and use the same row indices across all 14 representation points. 

To characterize the native representation geometry, effective rank is computed directly from the extracted forecast-slot representations, before they are standardized using training-set statistics. Probe parameters, predictions, and forecast targets do not enter the calculation.

As a stability check, we recompute effective rank over $200$ random \(80\%\) subsamples drawn without replacement. As shown in Figure~\ref{fig:effective_rank_stability}, the layerwise profiles remain stable, and the peak effective-rank layer is preserved across all 200 subsamples for each dataset.

\begin{figure}[t]
    \centering
    \includegraphics[width=\linewidth]{figures/erank_stability_3x3__fslot__probe_input__train.png}
    \caption{
    Stability of layerwise effective-rank estimates under subsampling.
    Solid curves show effective rank computed from the full sample of
    \(N=4096\) forecast-slot states, and shaded regions show the
    $2.5\text{th} - 97.5\text{th}$ percentile range across 200 random \(80\%\) subsamples
    (\(m=3277\)) drawn without replacement.
    The layerwise profiles are stable across subsamples, and the peak
    effective-rank representation point is preserved in all 200
    subsamples for each dataset.
    }
    \label{fig:effective_rank_stability}
\end{figure}

\subsection{Centered Kernel Alignment}
\label{apndx:cka}

Centered kernel alignment (CKA) provides a way to compare representations from different parts of a neural network \citep{kornblith2019similarityneuralnetworkrepresentations, cortes2024algorithmslearningkernelsbased}. For two representation matrices \(X\in\mathbb{R}^{N\times d_x}\) and \(Y\in\mathbb{R}^{N\times d_y}\), where \(N\) is the number of observations and \(d_x\) and \(d_y\) are the representation dimensions, we center each feature across the \(N\) observations:
\[
X_c = X-\bar{X}, \qquad Y_c = Y-\bar{Y},
\]
and then compute
\[
\mathrm{CKA}(X,Y)
=
\frac{\lVert X_c^\top Y_c\rVert_F^2}
{\lVert X_c^\top X_c\rVert_F
 \lVert Y_c^\top Y_c\rVert_F},
\]
where \(\|\cdot\|_F\) denotes the Frobenius norm.

Each representation point contains \(K=4\) forecast-slot states of dimension \(d=768\) per forecasting window, which we reshape before computing CKA:
\[
(n,K,d)\rightarrow(nK,d),
\]
where $n=262$ is the number of forecasting windows in the ID test set for a given dataset. 

We evaluate CKA for every pair of the 14 representation points (\(\mathrm{Emb}\), \(L1 - L12\), and \(L12 + \mathrm{RMS}\)), which gives a symmetric \(14\times14\) layer-similarity matrix with a unit diagonal. Probe parameters, predictions, and forecast targets do not enter the CKA calculation.

Figure~\ref{fig:additional-cka} shows the CKA matrices for the tested datasets. The representations exhibit one or more blocks of high similarity, but the location and sharpness of these structures vary across datasets and do not consistently align with the validation-selected saturation points.

\section{Data and Evaluation}

\subsection{Datasets}
\label{apndx:datasets}

We use four hourly datasets presented in the Chronos-2 pretraining corpus
\citep{ansari2025chronos2univariateuniversalforecasting}:

\paragraph{M4 Hourly.}
M4 Hourly consists of 414 anonymized real-world hourly time series from the M4 forecasting competition, obtained from various domains \citep{MAKRIDAKIS202054}. The series range from 748 to 1,008 observations in length.

\paragraph{Electricity.}
Electricity contains hourly electricity-load measurements from 321 clients in the Monash Time Series Forecasting Archive \citep{godahewa2021monashtimeseriesforecasting}. Each series contains 26,304 hourly observations.

\paragraph{Uber TLC.}
Uber TLC contains hourly Uber pickup counts across locations in New York City, derived from data released by the New York City Taxi and Limousine Commission \citep{fivethirtyeight_uber_tlc}. The dataset contains 262 series, each spanning 4,344 hours.

\paragraph{Wind Farms.}
Wind Farms contains hourly power-production measurements from 337 Australian wind farms \citep{godahewa2021monashtimeseriesforecasting}. The series vary in length from 1,715 to 8,784 observations. 
% About 17\% of values are missing, so after discarding windows with non-finite values, we retain 328 eligible series (Appendix \ref{apndx:data-splits}).

For experiments involving datasets outside the documented Chronos-2 pretraining corpus, we use:

\paragraph{BOOM.}
BOOM contains real-world observability metrics from Datadog for evaluating forecasting models in complex environments such as infrastructure, networking, databases, security, and applications \citep{cohen2025timedifferentobservabilityperspective}. We restrict the dataset to its 378 native-hourly metric queries.

\paragraph{SG Carpark.}
SG Carpark contains available parking-lot counts from 354 public carparks in Singapore. We use the 15-minute version provided through the TIME benchmark \citep{qiao2026itstimegenerationtime}.

\paragraph{Coastal T-S.}
Coastal T-S contains sea-water temperature, practical salinity, and relative-pressure measurements from 24 Australian coastal moorings. We use the hourly version distributed through TIME \citep{qiao2026itstimegenerationtime}.


\subsection{Preprocessing and rolling-origin splits}
\label{apndx:data-splits}

\paragraph{Rolling-origin splits.}
The datasets contain different numbers of time series with different lengths, so we construct examples using a rolling-origin split with non-overlapping forecast horizons. Within each series, the final forecast origin is assigned to test, the preceding origin to validation, and all earlier origins to training. For the PT-ID datasets, we set the train and validation/test budgets by the limiting dataset in each case: 1,394 training, 262 validation, and 262 test windows per dataset. PT-OOD evaluation uses all eligible test windows after preprocessing and filtering.

\paragraph{Dataset-specific preprocessing.}
For BOOM, we apply an additional pre-selection rule when choosing one variate from each query: a candidate is accepted only if at most \(20\%\) of its values are missing and it yields at least one valid forecasting window. For Coastal T-S, we keep temperature and salinity and discard relative pressure, which gives us 48 univariate series. For SG Carpark, 15-minute observations are aggregated to hourly resolution when at least three of four samples are available.  

\paragraph{Window construction and filtering.}
Windows with non-finite context/target values or near-constant contexts (\(\sigma<10^{-6}\)) are discarded. A series is retained only if at least three \(H\)-spaced windows remain: one for validation, one for test, and at least one for training. After preprocessing and filtering, 354 BOOM queries remain eligible. Wind Farms contains approximately \(17\%\) missing values, and 328 of the original 337 series remain eligible after filtering.

\section{Implementation Details}
\label{app:implementation}

This section collects implementation settings that are not naturally described alongside the corresponding methods. The optimizer, training schedule, and candidate weight-decay grid are fixed across experiments. Weight decay is the only training hyperparameter selected separately for each dataset and layer, always using validation loss.

\paragraph{Probe fitting.}
The shared forecast-slot probe is fit on the \(K\) forecast-slot states of the training windows. Each of the 768 representation dimensions is standardized using statistics computed from the training slot states at that layer, and the same statistics are applied at validation and test. The probe is a single
\(\mathbb{R}^{768} \rightarrow \mathbb{R}^{Q \times P}\) linear map shared across slots. We train it with AdamW at a learning rate of \(10^{-2}\) for 300 full-batch epochs without early stopping. Weight decay is applied to the weight matrix only, with the bias left unregularized, and is selected per layer from
\(\{10^{-5},10^{-4},10^{-3},10^{-2},10^{-1},0.3,1,3\}\)
using validation loss. Probes are fit under three initializations with seeds 0, 1, and 2. The backbone and training data are fixed across these runs, so probe initialization is the source of variation intentionally varied across fits.

\paragraph{Adapter fitting.}
The linear adapter uses the same optimizer, learning rate, training duration, and weight-decay grid as the probe. Weight decay is again applied only to the weight matrix and selected using validation loss. Each adapter is initialized to the identity and fit once from this fixed initialization. We therefore do not average adapter results over multiple initialization seeds, and adapter confidence intervals reflect test-set bootstrap uncertainty. Adapters are trained through the frozen final RMSNorm and forecasting head against the normalized forecast targets. All Chronos-2 parameters remain frozen, with only the adapter parameters updated.

\paragraph{Effective rank.}
Effective rank is computed on the training split from the \(K\) forecast-slot states stacked into an
\((N \cdot K) \times 768\) matrix. We use the representations before the per-layer standardization applied to the probe, so the resulting effective rank characterizes the representation itself. For each layer, we draw 4,096 forecast-slot states and report the mean effective rank over 200 subsamples containing 80\% of these states, sampled without replacement with seed 0.

\paragraph{Uncertainty and seeds.}
The global seed is 0 for NumPy, PyTorch, and scikit-learn, and window construction is deterministic given this seed, so all conditions are evaluated on identical windows. Confidence intervals use 5,000 cluster-bootstrap resamples with seed 0. Whole test series form the resampling unit for PT-ID datasets, while the corresponding parent cluster, namely carpark, station, or metric query, is used for PT-OOD datasets. A single resampling matrix is generated per dataset and reused across layers, conditions, and metrics to preserve paired comparisons.

\paragraph{Hardware and software.}
Experiments were run on a single NVIDIA A100-SXM4-40GB GPU with an AMD EPYC host CPU under SLURM. The software stack consists of Python 3.11.4, PyTorch 2.12.1 with CUDA 13.2 and cuDNN 9.21.1, and \texttt{chronos-forecasting} 2.3.1 for the \texttt{amazon/chronos-2} pipeline. We additionally use NumPy 1.26.4, scikit-learn 1.7.1, and Matplotlib 3.10.8. All computations use \texttt{float32}; we do not use \texttt{torch.compile}, mixed precision, or manual kernel fusion. The latency and memory benchmarking protocol is described separately in Appendix~\ref{app:latency-methodology}.

\paragraph{Reproducing the results.}
Each experimental stage has a separate entry point and writes the records consumed by subsequent stages. Representation extraction and probe fitting, including saturation-point selection, use
\texttt{run\_ptood\_probing\_ftok}; effective rank uses
\texttt{run\_spectral}; layer similarity uses
\texttt{run\_cka\_analysis}; and the frozen-head, adapter, and native conditions use
\texttt{run\_native\_head\_adapter}. Compression-cost tables and depth-selection rules are produced by
\texttt{run\_compression\_cost}, latency and memory measurements by
\texttt{run\_latency}, and paper figures by
\texttt{make\_id\_paper\_figures}. Representation extraction and model fitting are run on a GPU, while the tables, summary statistics, and figures are regenerated from the stored records on CPU. The aggregation stages verify the recomputed quantities against the stored records before producing the final outputs.

% \paragraph{Code availability.} Code and result records are available at \url{...}.

\section{Latency and Memory Methodology}
\label{app:latency-methodology}

All latency and memory measurements reported in this paper are obtained directly from the native and physically truncated models rather than inferred from parameter or FLOP counts. This appendix describes the benchmarking protocol used to obtain these measurements. The harness is implemented in \texttt{experiments/run\_latency.py} and evaluates the native model and each truncated variant under the same protocol.

\paragraph{Physical truncation and verification.}
The layerwise forecasting experiments in the main text obtain intermediate representations by running the full 12-block encoder and recording its hidden states with forward hooks. This produces the same representation at layer \(\ell\) as a model physically truncated after that layer, since later encoder blocks cannot influence \(h_\ell\). For latency and memory measurements, however, we explicitly construct the truncated model by removing encoder blocks \(\ell+1,\ldots,12\) and placing the trained linear adapter before the frozen final RMSNorm. Before benchmarking we perform two checks. First, at every measured depth, the physically truncated encoder must reproduce the corresponding states extracted from the full model. Second, at depth 12 with an identity-initialized adapter, the spliced path must reproduce the unmodified model's final quantile forecasts, which verifies the adapter and RMSNorm splice rather than only the block removal. The maximum absolute deviation was \texttt{0.0e+00} in both cases, so the truncated models are numerically identical to the hooked reference under this verification.

\paragraph{Hardware.}
Experiments were performed on a single NVIDIA A100-SXM4-40GB GPU with 39\,GiB of VRAM and compute capability 8.0, using driver version 580.178.04. The host machine used an AMD EPYC 7543 32-Core Processor, with two CPU cores allocated to the benchmarking job.

\paragraph{Software and precision.}
The benchmark used Python 3.11.4, PyTorch 2.12.1, CUDA 13.2, and \texttt{amazon/chronos-2}. The cuDNN version was 9.21.1 (\texttt{torch.backends.cudnn.version()} \(=\) 92101). All computations use \texttt{float32}, matching the precision used for the forecasting experiments in the main text. We do not use \texttt{torch.compile}, automatic mixed precision, or manual kernel fusion. TF32 matrix multiplication is disabled, while TF32 for cuDNN operations is enabled.

\paragraph{Workload.}
We use the same forecasting configuration as in the main experiments, with context length \(C=512\) and forecast horizon \(H=64\). This gives \(K=4\) forecast slots and an encoder sequence of 37 tokens, consisting of 32 context patches, one REG token, and four forecast slots. We benchmark batch sizes 1, 32, and 256 to cover both single-request and batched-throughput regimes. Inputs are synthetic random walks. Because the model has no data-dependent control flow, the computational path is unchanged across input values, so the benchmark does not depend on a particular dataset.

\paragraph{Timing protocol.}
Each configuration is first run for 20 warm-up iterations, which are discarded, followed by 100 timed repetitions. Since CUDA operations are asynchronous, \texttt{torch.cuda.synchronize()} is called before and after each timed region so that the reported latency includes GPU execution rather than only CPU-side dispatch. We report the median latency across repetitions as the primary statistic and the 95th percentile as a measure of tail latency.

We record two complementary latency measurements. \texttt{predict} measures \texttt{Chronos2Pipeline.predict\_quantiles} from host-resident inputs and therefore includes preprocessing and host-to-device transfer. In contrast, \texttt{encode} measures only the encoder forward pass on an input tensor already resident on the GPU. Speedups reported in the main text are based on \texttt{predict}, which provides the more conservative end-to-end comparison because its fixed preprocessing and transfer costs do not decrease with truncation.

\paragraph{Throughput.}
Throughput is reported in series per second and is computed as the batch size divided by the median \texttt{predict} latency. Larger batches amortize fixed per-call costs and make better use of the GPU, so throughput can increase substantially even though the latency of processing an entire batch also increases.

\paragraph{Memory.}
Peak GPU memory is measured with \texttt{torch.cuda.max\_memory\_allocated()} after resetting the peak-memory statistics with \texttt{torch.cuda.reset\_peak\_memory\_stats()}. We use PyTorch's allocated-memory measurement rather than \texttt{nvidia-smi}, whose reported reservation includes memory retained by the CUDA caching allocator and therefore does not directly reflect the memory released by truncation. For each depth, the model is loaded and truncated independently so that parameters belonging to removed blocks are released before measurement and allocations from a previous configuration do not affect the result.

\paragraph{Measured results.}
Table~\ref{tab:latency-full} reports the measured latency, throughput, and peak GPU memory across truncation depths and batch sizes. These quantities depend on truncation depth and batch size rather than on the forecasting dataset. Consequently, the dataset-specific speedups reported in the main text are obtained by evaluating this table at the truncation depth selected for each dataset, relative to the \texttt{Native} row at the same batch size.

\begin{table}[t]
    \centering
    \small
    \caption{
    Measured latency, throughput, and peak GPU memory across truncation depths.
    \texttt{Native} is the unmodified pretrained model; the \(L\ell\) rows are truncated at \(\ell\) blocks and carry the linear adapter, so \(L12\) is not the native model.
    \texttt{predict} measures the full prediction path from host-resident inputs, including preprocessing and host-to-device transfer, while \texttt{encode} measures the encoder forward pass on a device-resident tensor.
    }
    \label{tab:latency-full}
    \begin{tabular}{llrrrrrr}
        \toprule
        Batch & Depth & Blocks & predict (ms) & p95 (ms) & encode (ms) & series/s & peak (MiB) \\
        \midrule
        1 & \texttt{Native} & 12 & 18.45 & 18.75 & 17.65 & 54 & 466.3 \\
         & \(L1\) & 1 & 3.60 & 3.75 & 2.74 & 278 & 72.3 \\
         & \(L3\) & 3 & 6.34 & 6.49 & 5.47 & 158 & 144.5 \\
         & \(L6\) & 6 & 10.42 & 10.68 & 9.56 & 96 & 252.5 \\
         & \(L8\) & 8 & 13.10 & 13.43 & 12.23 & 76 & 324.5 \\
         & \(L10\) & 10 & 15.78 & 16.09 & 14.96 & 63 & 396.5 \\
         & \(L11\) & 11 & 17.17 & 17.47 & 16.38 & 58 & 432.5 \\
         & \(L12\) & 12 & 18.53 & 18.84 & 17.74 & 54 & 468.5 \\
        \midrule
        32 & \texttt{Native} & 12 & 26.08 & 26.33 & 24.24 & 1227 & 511.0 \\
         & \(L1\) & 1 & 5.33 & 5.41 & 3.47 & 6008 & 113.7 \\
         & \(L3\) & 3 & 9.13 & 9.30 & 7.27 & 3507 & 190.0 \\
         & \(L6\) & 6 & 14.81 & 15.04 & 12.96 & 2161 & 298.0 \\
         & \(L8\) & 8 & 18.59 & 18.84 & 16.74 & 1722 & 370.0 \\
         & \(L10\) & 10 & 22.39 & 22.67 & 20.54 & 1429 & 442.1 \\
         & \(L11\) & 11 & 24.31 & 24.58 & 22.44 & 1317 & 478.1 \\
         & \(L12\) & 12 & 26.21 & 26.45 & 24.36 & 1221 & 513.3 \\
        \midrule
        256 & \texttt{Native} & 12 & 178.07 & 178.54 & 168.51 & 1438 & 956.4 \\
         & \(L1\) & 1 & 27.76 & 28.12 & 18.29 & 9222 & 534.7 \\
         & \(L3\) & 3 & 55.19 & 55.52 & 45.72 & 4639 & 634.5 \\
         & \(L6\) & 6 & 96.39 & 96.73 & 86.83 & 2656 & 742.5 \\
         & \(L8\) & 8 & 123.80 & 124.08 & 114.25 & 2068 & 814.6 \\
         & \(L10\) & 10 & 151.20 & 151.53 & 141.65 & 1693 & 886.6 \\
         & \(L11\) & 11 & 164.94 & 165.29 & 155.40 & 1552 & 922.6 \\
         & \(L12\) & 12 & 178.70 & 179.07 & 169.13 & 1433 & 958.6 \\
        \bottomrule
    \end{tabular}
\end{table}

%%%%%%%%%%%%%%%%%%%%%%%%%%%%%%%%%%%%%%%%%%%%%%%%%%%%%%%%%%%%

\clearpage
\section*{NeurIPS Paper Checklist}

\begin{enumerate}

\item {\bf Claims}
    \item[] Question: Do the main claims made in the abstract and introduction accurately reflect the paper's contributions and scope?
    \item[] Answer: \answerYes{}
    \item[] Justification: The three claims in the abstract map onto Sections~\ref{sec:tunnel}--\ref{sec:model-setup-cost}. Layerwise saturation under the \(\varepsilon=5\%\) validation criterion on all seven datasets is shown in Section~\ref{sec:tunnel} and Table~\ref{tab:saturation-points}; the effective-rank and CKA analyses in Section~\ref{sec:geometry} and Appendices~\ref{apndx:geometry} and \ref{apndx:cka} support the claim that geometry shows layerwise structure but does not identify a universal transition; adapter recovery and its inconsistent cross-dataset transfer are shown in Sections~\ref{sec:dataset-specific-adapter} and \ref{sec:transfer-adapter}; and the quantitative compression numbers (\(29.4\%\) active parameters at \(L3\), \(\times 3.23\) speedup, \(34\%\) peak-memory reduction at batch size 256) come from Table~\ref{tab:compression-selection} and the measured benchmark in Table~\ref{tab:latency-full}. The scope is stated explicitly: one model (Chronos-2), seven hourly datasets, fixed \(C=512\) and \(H=64\), and a dataset-dependent rather than universal truncation depth.
    \item[] Guidelines:
    \begin{itemize}
        \item The answer \answerNA{} means that the abstract and introduction do not include the claims made in the paper.
        \item The abstract and/or introduction should clearly state the claims made, including the contributions made in the paper and important assumptions and limitations. A \answerNo{} or \answerNA{} answer to this question will not be perceived well by the reviewers. 
        \item The claims made should match theoretical and experimental results, and reflect how much the results can be expected to generalize to other settings. 
        \item It is fine to include aspirational goals as motivation as long as it is clear that these goals are not attained by the paper. 
    \end{itemize}

\item {\bf Limitations}
    \item[] Question: Does the paper discuss the limitations of the work performed by the authors?
    \item[] Answer: \answerYes{}
    \item[] Justification: Limitations are stated throughout the results and in Section~\ref{sec:model-setup-cost} and the conclusion: the saturation depth is dataset-dependent, neither the saturation nor the effective-rank criterion determines a universal truncation depth, cross-dataset adapter transfer is inconsistent (Section~\ref{sec:transfer-adapter}), and the Coastal T-S estimates are imprecise because that dataset provides only 24 bootstrap clusters. The conclusion further notes that the findings are established for a single model family at one frequency and horizon and identifies generalization across frequencies, horizons, and TSFM families as future work.
    \item[] Guidelines:
    \begin{itemize}
        \item The answer \answerNA{} means that the paper has no limitation while the answer \answerNo{} means that the paper has limitations, but those are not discussed in the paper. 
        \item The authors are encouraged to create a separate ``Limitations'' section in their paper.
        \item The paper should point out any strong assumptions and how robust the results are to violations of these assumptions (e.g., independence assumptions, noiseless settings, model well-specification, asymptotic approximations only holding locally). The authors should reflect on how these assumptions might be violated in practice and what the implications would be.
        \item The authors should reflect on the scope of the claims made, e.g., if the approach was only tested on a few datasets or with a few runs. In general, empirical results often depend on implicit assumptions, which should be articulated.
        \item The authors should reflect on the factors that influence the performance of the approach. For example, a facial recognition algorithm may perform poorly when image resolution is low or images are taken in low lighting. Or a speech-to-text system might not be used reliably to provide closed captions for online lectures because it fails to handle technical jargon.
        \item The authors should discuss the computational efficiency of the proposed algorithms and how they scale with dataset size.
        \item If applicable, the authors should discuss possible limitations of their approach to address problems of privacy and fairness.
        \item While the authors might fear that complete honesty about limitations might be used by reviewers as grounds for rejection, a worse outcome might be that reviewers discover limitations that aren't acknowledged in the paper. The authors should use their best judgment and recognize that individual actions in favor of transparency play an important role in developing norms that preserve the integrity of the community. Reviewers will be specifically instructed to not penalize honesty concerning limitations.
    \end{itemize}

\item {\bf Theory assumptions and proofs}
    \item[] Question: For each theoretical result, does the paper provide the full set of assumptions and a complete (and correct) proof?
    \item[] Answer: \answerNA{}
    \item[] Justification: The paper is entirely empirical and contains no theorems, formal claims, or proofs. Effective rank and CKA are stated as definitions (Appendices~\ref{apndx:effective-rank} and \ref{apndx:cka}) rather than as theoretical results.
    \item[] Guidelines:
    \begin{itemize}
        \item The answer \answerNA{} means that the paper does not include theoretical results. 
        \item All the theorems, formulas, and proofs in the paper should be numbered and cross-referenced.
        \item All assumptions should be clearly stated or referenced in the statement of any theorems.
        \item The proofs can either appear in the main paper or the supplemental material, but if they appear in the supplemental material, the authors are encouraged to provide a short proof sketch to provide intuition. 
        \item Inversely, any informal proof provided in the core of the paper should be complemented by formal proofs provided in appendix or supplemental material.
        \item Theorems and Lemmas that the proof relies upon should be properly referenced. 
    \end{itemize}

    \item {\bf Experimental result reproducibility}
    \item[] Question: Does the paper fully disclose all the information needed to reproduce the main experimental results of the paper to the extent that it affects the main claims and/or conclusions of the paper (regardless of whether the code and data are provided or not)?
    \item[] Answer: \answerYes{}
    \item[] Justification: Appendix~\ref{app:implementation} specifies the probe and adapter architectures, optimizer, learning rate, number of epochs, the weight-decay grid and the validation-based selection rule, the seeds, and the bootstrap protocol; Appendix~\ref{apndx:data-splits} specifies the rolling-origin split, window budgets, filtering rules, and dataset-specific preprocessing; Appendix~\ref{app:latency-methodology} specifies the truncation construction, its numerical verification, and the full latency and memory protocol. The model is the public \texttt{amazon/chronos-2} checkpoint used frozen through the public \texttt{chronos-forecasting} pipeline, and all seven datasets are public, so the reported experiments can be reproduced from the descriptions alone.
    \item[] Guidelines:
    \begin{itemize}
        \item The answer \answerNA{} means that the paper does not include experiments.
        \item If the paper includes experiments, a \answerNo{} answer to this question will not be perceived well by the reviewers: Making the paper reproducible is important, regardless of whether the code and data are provided or not.
        \item If the contribution is a dataset and\slash or model, the authors should describe the steps taken to make their results reproducible or verifiable. 
        \item Depending on the contribution, reproducibility can be accomplished in various ways. For example, if the contribution is a novel architecture, describing the architecture fully might suffice, or if the contribution is a specific model and empirical evaluation, it may be necessary to either make it possible for others to replicate the model with the same dataset, or provide access to the model. In general. releasing code and data is often one good way to accomplish this, but reproducibility can also be provided via detailed instructions for how to replicate the results, access to a hosted model (e.g., in the case of a large language model), releasing of a model checkpoint, or other means that are appropriate to the research performed.
        \item While NeurIPS does not require releasing code, the conference does require all submissions to provide some reasonable avenue for reproducibility, which may depend on the nature of the contribution. For example
        \begin{enumerate}
            \item If the contribution is primarily a new algorithm, the paper should make it clear how to reproduce that algorithm.
            \item If the contribution is primarily a new model architecture, the paper should describe the architecture clearly and fully.
            \item If the contribution is a new model (e.g., a large language model), then there should either be a way to access this model for reproducing the results or a way to reproduce the model (e.g., with an open-source dataset or instructions for how to construct the dataset).
            \item We recognize that reproducibility may be tricky in some cases, in which case authors are welcome to describe the particular way they provide for reproducibility. In the case of closed-source models, it may be that access to the model is limited in some way (e.g., to registered users), but it should be possible for other researchers to have some path to reproducing or verifying the results.
        \end{enumerate}
    \end{itemize}

\item {\bf Open access to data and code}
    \item[] Question: Does the paper provide open access to the data and code, with sufficient instructions to faithfully reproduce the main experimental results, as described in supplemental material?
    \item[] Answer: \answerNo{}
    \item[] Justification: Code is not attached to this submission. All datasets are public and are identified by their original sources and distribution channels (Appendix~\ref{apndx:datasets}), the model is the public \texttt{amazon/chronos-2} checkpoint, and Appendices~\ref{apndx:data-splits}--\ref{app:latency-methodology} document the preprocessing, split construction, hyperparameters, selection rules, and the per-stage entry points used to produce every reported number, which we intend to be sufficient for independent reimplementation. We plan to release the code and stored result records with the camera-ready version.
    % If an anonymized code release is attached to the submission, replace the answer above with
    % \answerYes{} and this justification with:
    % Code, configuration, and the stored result records needed to regenerate every table and figure
    % are provided in the supplementary material, together with the per-stage entry points listed in
    % Appendix~\ref{app:implementation}. All seven datasets are public and are obtained from the
    % sources named in Appendix~\ref{apndx:datasets}; the preprocessing and split construction are
    % implemented in the released code and documented in Appendix~\ref{apndx:data-splits}.
    % Remember to uncomment the Code availability paragraph in paper.tex as well.
    \item[] Guidelines:
    \begin{itemize}
        \item The answer \answerNA{} means that paper does not include experiments requiring code.
        \item Please see the NeurIPS code and data submission guidelines (\url{https://neurips.cc/public/guides/CodeSubmissionPolicy}) for more details.
        \item While we encourage the release of code and data, we understand that this might not be possible, so \answerNo{} is an acceptable answer. Papers cannot be rejected simply for not including code, unless this is central to the contribution (e.g., for a new open-source benchmark).
        \item The instructions should contain the exact command and environment needed to run to reproduce the results. See the NeurIPS code and data submission guidelines (\url{https://neurips.cc/public/guides/CodeSubmissionPolicy}) for more details.
        \item The authors should provide instructions on data access and preparation, including how to access the raw data, preprocessed data, intermediate data, and generated data, etc.
        \item The authors should provide scripts to reproduce all experimental results for the new proposed method and baselines. If only a subset of experiments are reproducible, they should state which ones are omitted from the script and why.
        \item At submission time, to preserve anonymity, the authors should release anonymized versions (if applicable).
        \item Providing as much information as possible in supplemental material (appended to the paper) is recommended, but including URLs to data and code is permitted.
    \end{itemize}

\item {\bf Experimental setting/details}
    \item[] Question: Does the paper specify all the training and test details (e.g., data splits, hyperparameters, how they were chosen, type of optimizer) necessary to understand the results?
    \item[] Answer: \answerYes{}
    \item[] Justification: Section~\ref{sec:model-setup} gives the forecasting configuration (\(C=512\), \(H=64\), \(P=16\), \(K=4\)) and the probe and adapter definitions, and Section~\ref{sec:data-eval} gives the datasets, metrics, and uncertainty protocol. Appendix~\ref{apndx:data-splits} details the rolling-origin train/validation/test construction and window budgets, and Appendix~\ref{app:implementation} details the optimizer (AdamW, learning rate \(10^{-2}\), 300 full-batch epochs, no early stopping), the weight-decay grid \(\{10^{-5},\dots,3\}\) selected per dataset and layer on validation loss, standardization statistics, and seeds. Weight decay is the only tuned hyperparameter and is never selected on test data.
    \item[] Guidelines:
    \begin{itemize}
        \item The answer \answerNA{} means that the paper does not include experiments.
        \item The experimental setting should be presented in the core of the paper to a level of detail that is necessary to appreciate the results and make sense of them.
        \item The full details can be provided either with the code, in appendix, or as supplemental material.
    \end{itemize}

\item {\bf Experiment statistical significance}
    \item[] Question: Does the paper report error bars suitably and correctly defined or other appropriate information about the statistical significance of the experiments?
    \item[] Answer: \answerYes{}
    \item[] Justification: All layerwise forecasting figures and the \(\Delta\)MASE column of Table~\ref{tab:compression-selection} report 95\% confidence intervals from 5{,}000 cluster-bootstrap resamples with seed 0, as described in Section~\ref{sec:data-eval} and Appendix~\ref{app:implementation}. The resampling unit is the whole test series for PT-ID datasets and the parent cluster (carpark, station, or metric query) for PT-OOD datasets, which accounts for correlation between windows drawn from the same series, and a single resampling matrix per dataset is reused across layers and conditions so that comparisons against native Chronos-2 are paired. Probe curves are additionally averaged over three initializations (seeds 0, 1, 2), whereas each adapter is fit once from its identity initialization, so adapter bands reflect test-set resampling only; this is stated in Appendix~\ref{app:implementation}. Effective-rank estimates are accompanied by a subsampling stability check over 200 subsamples (Figure~\ref{fig:effective_rank_stability}).
    \item[] Guidelines:
    \begin{itemize}
        \item The answer \answerNA{} means that the paper does not include experiments.
        \item The authors should answer \answerYes{} if the results are accompanied by error bars, confidence intervals, or statistical significance tests, at least for the experiments that support the main claims of the paper.
        \item The factors of variability that the error bars are capturing should be clearly stated (for example, train/test split, initialization, random drawing of some parameter, or overall run with given experimental conditions).
        \item The method for calculating the error bars should be explained (closed form formula, call to a library function, bootstrap, etc.)
        \item The assumptions made should be given (e.g., Normally distributed errors).
        \item It should be clear whether the error bar is the standard deviation or the standard error of the mean.
        \item It is OK to report 1-sigma error bars, but one should state it. The authors should preferably report a 2-sigma error bar than state that they have a 96\% CI, if the hypothesis of Normality of errors is not verified.
        \item For asymmetric distributions, the authors should be careful not to show in tables or figures symmetric error bars that would yield results that are out of range (e.g., negative error rates).
        \item If error bars are reported in tables or plots, the authors should explain in the text how they were calculated and reference the corresponding figures or tables in the text.
    \end{itemize}

\item {\bf Experiments compute resources}
    \item[] Question: For each experiment, does the paper provide sufficient information on the computer resources (type of compute workers, memory, time of execution) needed to reproduce the experiments?
    \item[] Answer: \answerYes{}
    \item[] Justification: Appendices~\ref{app:implementation} and \ref{app:latency-methodology} specify the compute worker and software stack: a single NVIDIA A100-SXM4-40GB GPU on an AMD EPYC 7543 host with two allocated CPU cores under SLURM, PyTorch 2.12.1 with CUDA 13.2, and \texttt{float32} throughout without \texttt{torch.compile} or mixed precision. Representation extraction, probe fitting, adapter fitting, and the latency benchmark run on that single GPU, and all tables, statistics, and figures are regenerated on CPU from stored records; the paper characterizes the hardware and the GPU/CPU split for each stage rather than itemizing per-run wall-clock times. Exploratory experiments not reported in the paper used the same single-GPU setup.
    \item[] Guidelines:
    \begin{itemize}
        \item The answer \answerNA{} means that the paper does not include experiments.
        \item The paper should indicate the type of compute workers CPU or GPU, internal cluster, or cloud provider, including relevant memory and storage.
        \item The paper should provide the amount of compute required for each of the individual experimental runs as well as estimate the total compute. 
        \item The paper should disclose whether the full research project required more compute than the experiments reported in the paper (e.g., preliminary or failed experiments that didn't make it into the paper). 
    \end{itemize}
    
\item {\bf Code of ethics}
    \item[] Question: Does the research conducted in the paper conform, in every respect, with the NeurIPS Code of Ethics \url{https://neurips.cc/public/EthicsGuidelines}?
    \item[] Answer: \answerYes{}
    \item[] Justification: The research uses a publicly released pretrained forecasting model and public benchmark datasets, involves no human subjects or personally identifying information (the traffic, parking, and observability datasets are aggregated counts and sensor measurements rather than individual-level records), and introduces no dual-use capability. The submission is anonymized.
    \item[] Guidelines:
    \begin{itemize}
        \item The answer \answerNA{} means that the authors have not reviewed the NeurIPS Code of Ethics.
        \item If the authors answer \answerNo, they should explain the special circumstances that require a deviation from the Code of Ethics.
        \item The authors should make sure to preserve anonymity (e.g., if there is a special consideration due to laws or regulations in their jurisdiction).
    \end{itemize}

\item {\bf Broader impacts}
    \item[] Question: Does the paper discuss both potential positive societal impacts and negative societal impacts of the work performed?
    \item[] Answer: \answerNo{}
    \item[] Justification: The paper does not include a dedicated broader-impacts discussion. The work is a methodological study of the inference cost of an already public forecasting model on public benchmarks; the main foreseeable positive impact is reduced inference latency, memory, and energy for deployed forecasting systems, and the main foreseeable negative impact is degraded forecast quality if a truncation depth is chosen on one dataset and reused on another, which the paper explicitly cautions against in Section~\ref{sec:model-setup-cost} and the conclusion.
    \item[] Guidelines:
    \begin{itemize}
        \item The answer \answerNA{} means that there is no societal impact of the work performed.
        \item If the authors answer \answerNA{} or \answerNo, they should explain why their work has no societal impact or why the paper does not address societal impact.
        \item Examples of negative societal impacts include potential malicious or unintended uses (e.g., disinformation, generating fake profiles, surveillance), fairness considerations (e.g., deployment of technologies that could make decisions that unfairly impact specific groups), privacy considerations, and security considerations.
        \item The conference expects that many papers will be foundational research and not tied to particular applications, let alone deployments. However, if there is a direct path to any negative applications, the authors should point it out. For example, it is legitimate to point out that an improvement in the quality of generative models could be used to generate Deepfakes for disinformation. On the other hand, it is not needed to point out that a generic algorithm for optimizing neural networks could enable people to train models that generate Deepfakes faster.
        \item The authors should consider possible harms that could arise when the technology is being used as intended and functioning correctly, harms that could arise when the technology is being used as intended but gives incorrect results, and harms following from (intentional or unintentional) misuse of the technology.
        \item If there are negative societal impacts, the authors could also discuss possible mitigation strategies (e.g., gated release of models, providing defenses in addition to attacks, mechanisms for monitoring misuse, mechanisms to monitor how a system learns from feedback over time, improving the efficiency and accessibility of ML).
    \end{itemize}
    
\item {\bf Safeguards}
    \item[] Question: Does the paper describe safeguards that have been put in place for responsible release of data or models that have a high risk for misuse (e.g., pre-trained language models, image generators, or scraped datasets)?
    \item[] Answer: \answerNA{}
    \item[] Justification: No data or models are released, and the assets used are already publicly available. Truncating the depth of an existing forecasting model creates no new risk of misuse, so no safeguards are required.
    \item[] Guidelines:
    \begin{itemize}
        \item The answer \answerNA{} means that the paper poses no such risks.
        \item Released models that have a high risk for misuse or dual-use should be released with necessary safeguards to allow for controlled use of the model, for example by requiring that users adhere to usage guidelines or restrictions to access the model or implementing safety filters. 
        \item Datasets that have been scraped from the Internet could pose safety risks. The authors should describe how they avoided releasing unsafe images.
        \item We recognize that providing effective safeguards is challenging, and many papers do not require this, but we encourage authors to take this into account and make a best faith effort.
    \end{itemize}

\item {\bf Licenses for existing assets}
    \item[] Question: Are the creators or original owners of assets (e.g., code, data, models), used in the paper, properly credited and are the license and terms of use explicitly mentioned and properly respected?
    \item[] Answer: \answerYes{}
    \item[] Justification: All assets are credited to their original publications and their distribution channels are named: the frozen \texttt{amazon/chronos-2} checkpoint accessed through \texttt{chronos-forecasting} 2.3.1 (Apache-2.0), the four pretraining-corpus datasets accessed through the \texttt{autogluon/chronos\_datasets} collection and cited to M4, the Monash archive, and the FiveThirtyEight release of the New York City Taxi and Limousine Commission data, BOOM (\texttt{Datadog/BOOM}, Apache-2.0), and SG Carpark and Coastal T-S as distributed by the TIME benchmark under CC BY-NC 4.0. Versions and software releases are recorded in Appendix~\ref{app:implementation}, and the two CC BY-NC 4.0 datasets are used only for non-commercial academic research, consistent with their terms.
    \item[] Guidelines:
    \begin{itemize}
        \item The answer \answerNA{} means that the paper does not use existing assets.
        \item The authors should cite the original paper that produced the code package or dataset.
        \item The authors should state which version of the asset is used and, if possible, include a URL.
        \item The name of the license (e.g., CC-BY 4.0) should be included for each asset.
        \item For scraped data from a particular source (e.g., website), the copyright and terms of service of that source should be provided.
        \item If assets are released, the license, copyright information, and terms of use in the package should be provided. For popular datasets, \url{paperswithcode.com/datasets} has curated licenses for some datasets. Their licensing guide can help determine the license of a dataset.
        \item For existing datasets that are re-packaged, both the original license and the license of the derived asset (if it has changed) should be provided.
        \item If this information is not available online, the authors are encouraged to reach out to the asset's creators.
    \end{itemize}

\item {\bf New assets}
    \item[] Question: Are new assets introduced in the paper well documented and is the documentation provided alongside the assets?
    \item[] Answer: \answerNA{}
    \item[] Justification: The paper introduces no new datasets, models, or other released assets. The trained linear probes and adapters are small artifacts used only to produce the reported measurements and are not distributed as assets.
    \item[] Guidelines:
    \begin{itemize}
        \item The answer \answerNA{} means that the paper does not release new assets.
        \item Researchers should communicate the details of the dataset\slash code\slash model as part of their submissions via structured templates. This includes details about training, license, limitations, etc. 
        \item The paper should discuss whether and how consent was obtained from people whose asset is used.
        \item At submission time, remember to anonymize your assets (if applicable). You can either create an anonymized URL or include an anonymized zip file.
    \end{itemize}

\item {\bf Crowdsourcing and research with human subjects}
    \item[] Question: For crowdsourcing experiments and research with human subjects, does the paper include the full text of instructions given to participants and screenshots, if applicable, as well as details about compensation (if any)? 
    \item[] Answer: \answerNA{}
    \item[] Justification: The paper involves no crowdsourcing and no research with human subjects; all experiments use existing public time-series datasets.
    \item[] Guidelines:
    \begin{itemize}
        \item The answer \answerNA{} means that the paper does not involve crowdsourcing nor research with human subjects.
        \item Including this information in the supplemental material is fine, but if the main contribution of the paper involves human subjects, then as much detail as possible should be included in the main paper. 
        \item According to the NeurIPS Code of Ethics, workers involved in data collection, curation, or other labor should be paid at least the minimum wage in the country of the data collector. 
    \end{itemize}

\item {\bf Institutional review board (IRB) approvals or equivalent for research with human subjects}
    \item[] Question: Does the paper describe potential risks incurred by study participants, whether such risks were disclosed to the subjects, and whether Institutional Review Board (IRB) approvals (or an equivalent approval/review based on the requirements of your country or institution) were obtained?
    \item[] Answer: \answerNA{}
    \item[] Justification: The paper involves no research with human subjects, so no IRB approval or equivalent review was required.
    \item[] Guidelines:
    \begin{itemize}
        \item The answer \answerNA{} means that the paper does not involve crowdsourcing nor research with human subjects.
        \item Depending on the country in which research is conducted, IRB approval (or equivalent) may be required for any human subjects research. If you obtained IRB approval, you should clearly state this in the paper. 
        \item We recognize that the procedures for this may vary significantly between institutions and locations, and we expect authors to adhere to the NeurIPS Code of Ethics and the guidelines for their institution. 
        \item For initial submissions, do not include any information that would break anonymity (if applicable), such as the institution conducting the review.
    \end{itemize}

\item {\bf Declaration of LLM usage}
    \item[] Question: Does the paper describe the usage of LLMs if it is an important, original, or non-standard component of the core methods in this research? Note that if the LLM is used only for writing, editing, or formatting purposes and does \emph{not} impact the core methodology, scientific rigor, or originality of the research, declaration is not required.
    %this research? 
    \item[] Answer: \answerNA{}
    \item[] Justification: No large language model is a component of the methodology. The studied model, Chronos-2, is a time-series foundation model, and every experimental component (linear probes, linear adapters, the frozen quantile head, and the representation diagnostics) is implemented directly.
    \item[] Guidelines:
    \begin{itemize}
        \item The answer \answerNA{} means that the core method development in this research does not involve LLMs as any important, original, or non-standard components.
        \item Please refer to our LLM policy in the NeurIPS handbook for what should or should not be described.
    \end{itemize}

\end{enumerate}



\end{document}